% marge
\usepackage[margin=1in]{geometry}
% langue
\usepackage[french]{babel}
\usepackage[T1]{fontenc}
\usepackage{microtype}
\usepackage{fontawesome5}
% police
% \usepackage{cmbright}
\usepackage[sfmath,light]{kpfonts}
\renewcommand{\familydefault}{\sfdefault}
% math utils
\usepackage{mathtools, amssymb}
\usepackage{autoaligne}
\usepackage{bm}
% couleurs
\usepackage[dvipsnames, svgnames]{xcolor}
% image
\usepackage{graphicx}
% dessin
\usepackage{tikz}
\usetikzlibrary{babel}
\usepackage{tkz-fct}
\usepackage{tkz-tab}
\usepackage{tkz-euclide}
\usepackage{pgfplots}
\pgfplotsset{compat=1.18}
\usepackage{tikzscale}
\usepackage{tcolorbox}
\tcbuselibrary{skins,xparse}
% plusieurs colonnes
\usepackage{multicol}
% soulignements différents
\usepackage{ulem}
% gestion des imports
\usepackage{import}
% tableaux
\usepackage{tabularray}
% personnaliser les listes
\usepackage{enumitem}
\setlist[enumerate]{label=\textbf{\alph*)}}
\usepackage{tasks}
\settasks{
    label-width=15 pt,
    label-format=\bfseries
}
% unités
\usepackage[locale=FR, group-minimum-digits=3]{siunitx}
%% mettre les labels sous forme de lettre en gras
%% nouveau paramètre pour mettre les listes sur plusieurs colonnes
\SetEnumitemKey{cols}{%
    itemsep = 1\itemsep,
    parsep = 1\parsep,
    before = \raggedcolumns\begin{multicols}{#1},
        after = \end{multicols}}
% exercice
% use-files -> xsim utilise mets chaque exercice dans un fichier séparé
% use-aux -> xsim utilise le fichier aux de latex au lien d'en créer un nouveau
\usepackage[use-files,use-aux]{xsim}
% XSIM TEMPLATE
\DeclareExerciseEnvironmentTemplate{math-exam}
{%
{\colorbox{black}{\color{white}\bfseries\large\GetExerciseProperty{counter}}~}
\marginpar{\quad/\GetExerciseProperty{points}}
}%
{\bigbreak}

% Mes environnements
\DeclareExerciseType{connaitre}{
    exercise-env = connaitre ,
    solution-env = solconnaitre ,
    exercise-name = \XSIMtranslate{question} ,
    exercises-name = \XSIMtranslate{questions} ,
    solution-name = \XSIMtranslate{solution} ,
    solutions-name = \XSIMtranslate{solutions} ,
    exercise-template = math-exam ,
    solution-template = math-exam ,
    exercise-heading = \subsection* ,
    solution-heading = \subsection*,
    counter = exercise
}
\DeclareExerciseType{appliquer}{
    exercise-env = appliquer ,
    solution-env = solappliquer ,
    exercise-name = \XSIMtranslate{question} ,
    exercises-name = \XSIMtranslate{questions} ,
    solution-name = \XSIMtranslate{solution} ,
    solutions-name = \XSIMtranslate{solutions} ,
    exercise-template = math-exam ,
    solution-template = math-exam ,
    exercise-heading = \subsection* ,
    solution-heading = \subsection*,
    counter = exercise
}
\DeclareExerciseType{transferer}{
    exercise-env = transferer ,
    solution-env = soltransferer ,
    exercise-name = \XSIMtranslate{question} ,
    exercises-name = \XSIMtranslate{questions} ,
    solution-name = \XSIMtranslate{solution} ,
    solutions-name = \XSIMtranslate{solutions} ,
    exercise-template = math-exam ,
    solution-template = math-exam ,
    exercise-heading = \subsection* ,
    solution-heading = \subsection*,
    counter = exercise
}
\xsimsetup{
    path=xsim,% stocker les fichiers auxiliaires dans ce dossier
    % solution/print=true, % afficher les solutions en dessous des questions
    blank/style=\dotuline{#1},
    exercise/template=math-exam,
    solution/template=math-exam,
    exercise/name=Question,
    % solconnaitre/print=true,
    % solappliquer/print=true,
    % soltransferer/print=true,
}

\newcommand{\systeme}[2]
{%
    $\left\{\begin{aligned}
        #1 \\
        #2
    \end{aligned}\right.$
}
